\section{Introduction and source problem}\label{sec:introduction}
Krylov constructions convert unitary evolution into motion along a basis
generated by successive applications of the dynamical generator. Operator
complexity measures the mean position in an operator Krylov chain, while
spread complexity uses a state-vector chain
\cite{Parker2019,Balasubramanian2022}. A density matrix can participate in
either construction: it is itself an operator, and it also specifies a pure
state after purification. Comparing these descriptions requires control of
both the seed normalization and the generator. Earlier work on density-matrix
operators already distinguishes their Krylov complexity from the spread
complexity of a pure state \cite{Caputa2024}.

Das and Mori \cite{DasMori2026} formulated this comparison for two unitary
purification schemes. The time-independent scheme evolves only the physical
system, while the time-dependent scheme evolves the ancilla by the conjugate
unitary. Their Eq.~(4) proposes
\begin{equation}\label{eq:source-hierarchy}
 \CS\bigl(\Psi_\rho^{U^*}(t)\bigr)
 \ \le\ \CK\bigl(\rho(t)\bigr)
 \ \le\ \CK\bigl(\Psi_\rho^I(t)\bigr).
\end{equation}
Here the rightmost quantity is the operator complexity of the rank-one
density matrix of the purification. It is not the spread complexity of the
purified vector. We retain the published source's conventions throughout.

The source also proposes two different kinds of statement. Its second
conjecture describes the ratio of the time-dependent purification's spread
complexity to the mixed-state operator complexity as approximately constant
relative to its temporal mean, with a value depending on purity. This is
qualitative: it specifies neither a uniform tolerance nor an exact function
of purity. The third conjecture is the precise factor-two inequality between
operator and spread complexities of the same purification. The distinction
between these statements is central to the analysis below.

We give separate exact qutrit witnesses against the two sides of
Eq.~\eqref{eq:source-hierarchy}. The source's qubit hierarchy remains valid;
its explicit three-atom reduction proves that neither inequality can fail in
a smaller physical dimension. For temporal ratios, we state and rule out
three universal formulations: a finite uniform recurrence-period coefficient
of variation for the main ratio, the corresponding bound for its reciprocal,
and finite purity-only upper control of the main ratio or its mean. The
two concentration questions are logically distinct, so they require different
families. All counterexamples use strictly positive density matrices at every
allowed parameter value.

The factor-two inequality has a different status. Tensor-product
superadditivity, stated by Murugan and van Zyl
\cite[Sec.~2, Eq.~(2.4)]{MuruganVanZyl2026}, implies it after vectorizing a
rank-one projector. We give the needed finite-dimensional proof in terms of
nested polynomial subspaces and apply it separately to the two purification
generators. This known implication is context for the failed mixed-state
comparisons, rather than a new headline inequality. The early-time
pure-density relation also has antecedents in Ref.~\cite{Caputa2024}.

Our focus is finite-dimensional, time-independent unitary dynamics. The
results do not assess nonunitary examples or typical-state limits. A final fresh
public-prior-art sweep before release found no matching
public prior result; absolute priority
is not claimed.

Section~\ref{sec:framework} fixes the definitions and spectral representation.
Sections~\ref{sec:hierarchy}--\ref{sec:ratios} contain the comparison results
and their proofs. Section~\ref{sec:survivors} gives restricted survivors and
the spectral mechanism. Exact finite certificates and technical estimates
are collected in the appendices.
